\section{NN-S013: zero avoidance, inversion, and necessary cancellation in generators}

\paragraph{Statement.}
Use the circuit models of NN-S012, and give multi-output circuits a
polynomial number of output bits. The following refinements hold.
\begin{enumerate}
\item Every map given by a circuit over $G(\mathbb F_2)$ that avoids
      supported zeros at every primitive gate can be inverted in
      deterministic polynomial time on its image, given the circuit.
      The algorithm also decides whether a proposed output has a preimage.
\item Ordinary one-way functions exist if and only if there are one-way
      functions with uniform polynomial-size circuits over
      $G(\mathbb F_4)$ avoiding supported zeros at every primitive gate.
      They also exist if and only if there are such functions over
      $G(\mathbb F_2)$ avoiding supported zeros only at the outputs.
      These equivalences retain the usual uniform or nonuniform
      adversary convention on both sides.
\item Let a polynomial-size circuit over $G(R)$, for a fixed finite
      nonzero ring $R$, have $n$ Boolean inputs and at least $n+1$ Boolean amplitude
      outputs. If its amplitude map is a pseudorandom generator secure
      against nonuniform polynomial-size distinguishers, then among
      its first $n+1$ outputs some output takes the supported zero on
      at least $\Omega(2^n/n^2)$ Boolean inputs, for all sufficiently
      large $n$. The selected output may depend on $n$.
\end{enumerate}
The third conclusion concerns the final gates themselves, not merely
some unspecified internal cancellation. There is also an unconditional
quantitative inequality, proved below, for every such map whether or
not it is pseudorandom. None of the assertions uses A or M.

The imported classical results are Goldreich--Izsak,
\emph{Monotone Circuits: One-Way Functions versus Pseudorandom
Generators}, Theory of Computing 8 (2012), 231--238:
Theorem 1.1, printed p.~232, and its proof in Section 3;
Theorem 1.2, printed p.~232, with the one-coordinate/two-coordinate
distinguishers constructed in Section 4.2, pp.~236--237.
Their Theorem 3.1 and proof, p.~234, provide the precise published
slice-circuit result attributed there to Berkowitz's 1982 unpublished
manuscript. No unlocated theorem number in that manuscript is assumed.

\paragraph{Proof of efficient inversion.}
The effective deletion in NN-S012 turns the promised circuit into
multiplication gates, constants and input variables. Associate to each
nonzero gate the subset of variables in its conjunction; associate a
separate zero marker to a gate that is identically $\tau$.
Multiplication takes unions unless one factor has the zero marker.
The constant one corresponds to the empty subset. These sets are
computed with polynomially many bit-set operations.

For a proposed output $y$, first reject if it assigns one to a constant
zero output or zero to a constant one output. Let $A$ be the union of
all conjunction subsets belonging to outputs where $y_i=1$. Set
$x_j=1$ for $j\in A$ and $x_j=0$ otherwise. Every proposed preimage
must have all bits in $A$ set to one, so this is the smallest possible
preimage in the coordinatewise order. It satisfies every output that
is required to be one. If a remaining nonconstant output is required
to be zero, with conjunction subset $S_i$, then it vanishes at this
$x$ exactly when $S_i\not\subseteq A$. If instead $S_i\subseteq A$,
every possible preimage would make that output one, so no preimage
exists. Testing these inclusions therefore either returns a correct
preimage or proves emptiness. The entire computation is polynomial
in the circuit and input/output sizes. For a uniform family the
circuit is itself polynomial-time constructible; for nonuniform
families it can be supplied as polynomial advice. Thus this family
cannot furnish one-wayness in the corresponding convention.

\paragraph{Proof of the one-way-function equivalences.}
By Goldreich--Izsak Theorem 1.1, existence of one-way functions implies
existence of such functions computed by uniform polynomial-size
monotone Boolean circuits. Replace their AND and OR gates by the
explicit every-gate zero-free $G(\mathbb F_4)$ gadgets of NN-S012.
This changes neither their Boolean input/output map nor their
inversion problem and only multiplies circuit size by a constant.
Alternatively use the output-only $G(\mathbb F_2)$ gadgets of that
finding. Their internal supported zeros are allowed in this second
model. Both conversions are uniform. Conversely, any alleged
one-way map in either finite-ring model is an ordinary efficiently
computable Boolean map by NN-S012's finite lookup-table simulation.
The same adversary and input distribution give ordinary one-wayness.
This proves both equivalences. The conclusion is conditional only
on existence of ordinary one-way functions; it does not assert that
they exist or that E preserves any stronger exponential security.

\paragraph{Proof of the necessary output-cancellation bound.}
Retain the first $n+1$ outputs. Let $F(x)\in\{0,1\}^{n+1}$ be their
amplitude vector and $S(x)\in\{0,1\}^{n+1}$ their support vector.
The latter is a monotone map, regardless of any internal cancellations.
For each output $i$ put
\[
 B_i=\{x:C_i(x)=e\}.
\]
Since the amplitude outputs are Boolean, the only possible split
values are $\tau,e,1$. Hence $F_i(x)\ne S_i(x)$ if and only if
$x\in B_i$.

Goldreich--Izsak's quantitative Theorem 1.2 and its explicit proof
give an absolute $c>0$, for every sufficiently large $n$, a set
$J\subseteq\{1,\ldots,n+1\}$ of size one or two, and a monotone
Boolean test $D$ depending only on those coordinates, such that
\[
 |\mathbb E D(S(U_n))-\mathbb E D(U_{n+1})|\ge c/n^2.
\]
The test is either a coordinate or the AND of two coordinates.
Couple $F(U_n)$ and $S(U_n)$ using the same input. Their $D$-values
agree outside $\bigcup_{i\in J}B_i$. The triangle inequality gives
the unconditional estimate
\[
 |\mathbb E D(F(U_n))-\mathbb E D(U_{n+1})|
 \ge \frac{c}{n^2}
       -\Pr\left[U_n\in\bigcup_{i\in J}B_i\right].       \tag{S013.1}
\]
For a pseudorandom generator the left side is negligible, since the
chosen tests form a family of constant-size nonuniform circuits.
Consequently the probability on the right is at least
$c/n^2-\operatorname{negl}(n)$. The union bound and $|J|\le2$
give, for at least one selected output,
\[
 \Pr[C_i(U_n)=e]\ge c/(2n^2)-\operatorname{negl}(n).
\]
Multiplication by $2^n$ proves the stated cardinality bound.
Truncating a longer pseudorandom output to its first $n+1$ bits
preserves security, so the argument applies to every positive stretch.
The quantitative statement is deliberately made with nonuniform
security; (S013.1) itself has no security or uniformity hypothesis.

\paragraph{Slice elimination and the complexity gate.}
Goldreich--Izsak Theorem 3.1 gives a polynomial-time construction of
a monotone circuit for a slice function from any Boolean circuit for
it. Combining that theorem with NN-S012 yields polynomial-cost
elimination of supported zeros at the output over $G(\mathbb F_2)$,
and at every primitive gate over $G(\mathbb F_4)$. It does not give
every-gate elimination over $G(\mathbb F_2)$: even the one-slice
extension of two-variable OR is OR itself and violates the conjunction
classification. Thus the slice theorem, the one-way-function theorem,
and the pseudorandom-generator obstruction have a consistent precise
translation, but use different notions of zero avoidance.

E supplies the sparse preimage-counting algorithms of NN-C002 and
NN-C011, and the instance-dependent isolating construction NN-C019.
Those theorems retain their sparsity and witness-size parameters.
The polynomial monotone circuits furnished by the classical one-way
theorem need not have linearly many bounded-arity factors in their
original input dimension, and a subexponential bound would in any
case not be polynomial inversion. No conclusion that E destroys
one-way functions is drawn. The cancellation lower bound (S013.1)
is a distributional necessary condition for pseudorandomness, not a
general Boolean or algebraic circuit-size lower bound.

\paragraph{Geometric formulation and exact dictionary.}
On the split Boolean finite scheme $X$ of NN-S012, an output fibre
is a clopen union of Boolean points. The inversion proof computes
emptiness and a smallest point of fibres of maps whose coordinate
pullbacks are conjunction monomials. A Boolean amplitude map has
a corresponding morphism between the split Boolean schemes, while
the support map retains an additional clopen tag at each coordinate.
The locus $B_i$ is the clopen intersection where that tag is present
and the amplitude section is zero. Uniform input measure is the
normalized counting measure on the $2^n$ rational points. Thus
(S013.1) is equivalently a lower bound on the rank of the finite
coordinate algebra of one such $B_i$, divided by $2^n$.

Finite-free trace and its base change are exactly EGA IV$_4$,
18.2.1--18.2.2, original printed p.~111; the finite-fibre sheaf trace
is SGA 4, XVII, 6.2.3. Every such finite rational-point scheme has
Frobenius acting as the identity on its degree-zero cohomology.
That purity does not construct its potentially exponentially large
point-indexed basis from the circuit. The missing effectiveness is
the same distinction as NN-C006. The finite extension from
$\mathbb F_2$ to $\mathbb F_4$ changes the available zero-free circuit
presentations; it leaves the underlying Boolean finite set and its
counting distribution unchanged. Neither the trace identity nor
the sign-fixed skeleton supplies a polynomial inverter for the
monotone maps allowed over $\mathbb F_4$.
\hfill$\square$
